%!TeX program = xelatex
\documentclass[12pt,hyperref,a4paper,UTF8]{ctexart}
\usepackage{UIBEReport}

%%-------------------------------正文开始---------------------------%%
\begin{document}

%%-----------------------期刊论文信息设置--------------------%%
\title{\songti\bfseries\zihao{2} 中国企业的人工智能投入 —— 来自招聘大数据的发现\thanks{作者简介：陈沫樵，中国金融学院鸿儒2401班硕士研究生，学号：202442087。}}
\author{\songti\zihao{4} 陈沫樵}
\date{}

\maketitle

%%------------------期刊摘要与关键词-------------%%
\begin{abstract}
本文主要介绍本 LaTeX 论文模板的排版格式与使用说明。本模板适用于期末小论文或学术报告撰写，支持标准的双语参考文献自动分类与排序，采用宋体字形和 22 磅固定行距，已取消独立封面与目录，以满足学术期刊的流式排版规范。

{\heiti 关键词：}论文模板；排版格式；双语文献；宋体

{\heiti JEL 分类号：}A10, B20, C30
\end{abstract}

\section{模板说明}
本模板主要适用于一些课程的平时论文以及期末论文\footnote{这是符合《金融类专业学位论文写作规范》的按页独立编号的脚注示例。}，默认页边距为2.5cm，中文宋体，英文Times New Roman，字号为12pt（小四），行距为固定值22磅。

编译方式：\verb|xelatex -> biber -> xelatex*2|


默认模板文件由以下四部分组成：
\begin{itemize}
    \item \texttt{main.tex} 主文件
    \item \texttt{reference.bib} 参考文献，使用bibtex
    \item \texttt{UIBEReport.sty} 文档格式控制，包括一些基础的设置，如页眉、标题、姓名等
    \item \texttt{figures} 放置图片的文件夹
\end{itemize}

第一次使用时需前往\texttt{UIBEReport.sty} 对标题、姓名、学号、院所、页眉等进行设置，设置完后即可一劳永逸，封面logo亦可替换

默认带有封面页以及目录页，页码从目录页开始

\section{一些插入功能}
\subsection{插入公式}
行内公式$v-\varepsilon+\phi=2$。

插入行间公式如\autoref{Euler}：
\begin{equation}
    v-\varepsilon+\phi=2
    \label{Euler}
\end{equation}

\subsection{插入图片}
图片示例如\autoref{UIBE}所示，注意这里使用了\verb|~\autoref{}|命令，也就是会自动生成“图”“式”等前缀，无需手动输入。

\begin{figure}[!htbp]
    \centering
    \begin{minipage}{0.5\textwidth}
        \centering
        \framebox[\textwidth]{\rule{0pt}{3cm} 图片占位符 (Placeholder)}
    \end{minipage}
    \caption{图片排版示例}
    \label{UIBE}
\end{figure}

插入上面图片的代码（实际写作时可替换为具体的图片路径，例如 \texttt{figures/my\_image.png}）：

\begin{verbatim}
    \begin{figure}[!htbp]
        \centering
        \includegraphics[width =.5\textwidth]{figures/my_image.png}
        \caption{图片排版示例}
        \label{UIBE}
    \end{figure}
\end{verbatim}

\subsection{插入文本框}
本模板定义了一个圆角灰底的文本框，使用简化命令\verb|\tbox{}|即可，如果你不喜欢，可以前往 \texttt{UIBEReport.sty}对其进行修改。

\tbox{
    这是一个圆角灰底的文本框
}

\subsection{插入表格}
本模板文件如\autoref{doc}所示。根据规范，表格采用三线表，字号为小五，且标题置于表格上方。在经济学写作中，常用 \texttt{threeparttable} 宏包来排版带有底部注释的三线表。

\begin{table}[!htbp]
    \centering
    \begin{threeparttable}
    \caption{本模板文件组成}
    \label{doc}
    \begin{tabular*}{\textwidth}{@{\extracolsep{\fill}}ll}
    \toprule
        文件名 & 说明 \\
    \midrule
        \texttt{main.tex}  & 主文件 \\
        \texttt{reference.bib} & 参考文献 \\
        \texttt{UIBEReport.sty}  & 文档格式控制\\
        \texttt{figures}  & 图片文件夹 \\
    \bottomrule
    \end{tabular*}
    \begin{tablenotes}
        \item 注：表中列出了本模板的核心文件组成，格式符合对外经济贸易大学学术与学术期刊的通行规范。
    \end{tablenotes}
    \end{threeparttable}
\end{table}

%\section{定理环境}
%\begin{Theorem}
%\end{Theorem}
%
%\begin{Lemma}
%\end{Lemma}
%
%\begin{Corollary}
%\end{Corollary}
%
%\begin{Proposition}
%\end{Proposition}
%
%\begin{Definition}
%\end{Definition}
%
%\begin{Example}
%\end{Example}
%
%\begin{proof}
%\end{proof}

\subsection{插入参考文献}
直接使用\verb|\parencite{}|（用于括号引用）或\verb|\textcite{}|（用于叙述引用）即可。本模板采用 APA 参考文献格式，且自动按照中英文分类：其中中文在前（按拼音首字母排序），英文在后（按英文字母排序）。

例如：

   此处括号引用了中文文献\parencite{xu2017}和\parencite{wang2014}，叙述引用了中文文献\textcite{xu2017}。此处括号引用了英文文献\parencite{todaro1969}和\parencite{banister1990}，叙述引用了英文文献\textcite{todaro1969}。

引用过的文献会自动出现在参考文献中。

\section{写在最后}
\subsection{发布地址}
\begin{itemize}
    \item Github: \url{https://github.com/MochiaoChen/UIBE_temp}
\end{itemize}

%%----------- 参考文献 -------------------%%
%在reference.bib文件中填写参考文献，此处自动生成

\reference


\end{document}
